\documentclass[letterpaper]{article} % DO NOT CHANGE THIS
\usepackage[submission]{aaai2027} % DO NOT CHANGE THIS
\usepackage[hyphens]{url} % DO NOT CHANGE THIS
\usepackage{graphicx} % DO NOT CHANGE THIS
\urlstyle{rm} % DO NOT CHANGE THIS
\def\UrlFont{\rm} % DO NOT CHANGE THIS
\usepackage{natbib} % DO NOT CHANGE THIS
\usepackage{caption} % DO NOT CHANGE THIS
\usepackage{booktabs}
\frenchspacing % DO NOT CHANGE THIS

\pdfinfo{
/TemplateVersion (2027.1)
}

\setcounter{secnumdepth}{0}
\newcommand{\drafttodo}[1]{\textbf{[TODO: #1]}}

\title{TODO: Working Title for AlemDICE}
\author{Anonymous Submission}
\affiliations{}

\begin{document}

\maketitle

\begin{abstract}
% Abstract version 1: evidence-focused framing, retained for comparison.
% Long-horizon multi-agent language-model systems need mechanisms for allocating
% work, maintaining coherent roles, and scaling beyond small teams without relying
% on a privileged central planner.  Before such mechanisms can be evaluated,
% however, their execution and accounting must be reproducible under concurrent,
% failure-prone model calls.  We present the current  research substrate,
% which combines the long-horizon Alem coordination environment with a Contract
% Net task market, conserved CU credits, award leases, verified settlement,
% reauction, deterministic replay, and detailed provenance and usage logging.
% The implementation also supports concurrent OpenAI Responses API decisions and
% interfaces with Alem's language-agent action space.  Eight-agent deterministic
% runs validate the market state machine and exact treasury conservation, while a
% separate one-tick hosted-model check validates concurrent structured decisions,
% action parsing, and token accounting.  These checks do not yet establish a
% behavioral benefit: the deterministic policy ignored CU, the existing market
% comparison contains privileged-state and interface confounds, and a 100-tick
% horizon produces little progression in Alem.  We therefore make no claim that
% Contract Net or CU improves coordination over language-mediated agents.
% \drafttodo{Insert the matched language-only versus language-plus-market pilot result.}
% This foundation is intended to support controlled studies of teams with 32 or
% more agents, decentralized auctions for tasks and roles, and long-horizon role
% coherence.
% Abstract version 2: vision-led framing.
We study emergent collective behavior in networked AI agents under decentralized
control, with an emphasis on market-based coordination targeting teams of up to
32 large-language-model and vision-language-model agents.  Traditional
distributed-AI systems often combine centralized orchestration with
task-specific policies, limiting transfer and making coordination brittle as
missions or teams change.  Large pretrained models can instead interpret
unstructured commands and adapt allocation and decision policies to unseen
problems without task-specific retraining.  We investigate how these
capabilities can be
coupled to decentralized Contract Net auctions and conserved CU credits for
task and role allocation.  Our experiments will compare language-only,
static-allocation, and language-plus-market coordination while scaling team
size, interaction horizon, and communication sparsity.  We will evaluate
mission success, allocation and settlement efficiency, communication cost,
recovery time, and role coherence.  To characterize robustness, we will inject
communication loss, network partitions, malfunctioning agents, and hostile or
misaligned behavior.  These studies address three central obstacles to
dependable collective intelligence: self-organized communication with bounded
network cost, resilience to individual-agent failure or compromise, and
long-term role alignment that preserves compliance with mission intent.  Our
goal is to establish when large pretrained models improve decentralized task
allocation, decision-making, and control, and which coordination and governance
mechanisms are required for scalable, robust agent collectives.
\end{abstract}

\section{Methods: Current Working Notes}

\subsection{Research Question and Current Claim Boundary}

The intended scientific comparison asks whether decentralized task and role
allocation through a Contract Net market and CU credits improves coordination
relative to Alem's language-mediated agents.  The current implementation
supports running that comparison, but the corrected matched experiment has not
yet completed.  Our evidence currently supports only the following claims:

\begin{itemize}
    \item The auction, CU ledger, award-lease, verification, settlement,
    reauction, persistence, replay, and logging mechanisms execute end to end.
    \item Concurrent hosted-model decisions, structured outputs, action parsing,
    and token accounting work for eight agents in a one-tick integration check.
    \item Existing deterministic runs do not test whether CU changes language-
    model behavior and do not provide a valid comparison with Alem's canonical
    language agents.
\end{itemize}

Scaling beyond the current tests, decentralized role auctions, and measurement
of role coherence are research targets rather than completed contributions.
Sparse peer-to-peer communication and process-failure experiments are likewise
not part of the current evidence.

\subsection{Alem Environment and Reproduction Boundary}

We build on Alem, a procedurally generated, long-horizon multi-agent survival
and crafting environment with explicit coordination tasks
\citep{tessera2026alem}.  The repository is pinned to upstream commit
\texttt{b1344e46}.  Environment mechanics are
inherited; AlemDICE changes the experiment and coordination infrastructure, not
the underlying world.  The pinned revision follows upstream current main and is
not a bit-identical reproduction of the public \texttt{v0.1.0} paper release.

The canonical language-agent condition uses three agents with warrior, forager,
and miner roles, text observations, broadcast communication, private scratchpad
memory, and progressive disclosure of game rules.  Agents act synchronously in
the same world while receiving role-specific and locally rendered observations.
Soft specialization is enabled and shared reward is disabled.  Alem evaluates
ordinary and coordination achievements separately across Easy, Medium, and Hard
coordination settings.

The audit identified an Easy-setting discrepancy that must be resolved before
the matched study: the intended Easy setting implies 70\% non-specialist
efficiency, whereas the released LLM configuration explicitly sets it to 20\%.
All treatments in the corrected experiment will use the same documented value.

\subsection{Language-Agent Interface}

The upstream language condition uses homogeneous zero-shot instances of one
model.  Each decision includes a text observation, currently legal actions,
visible coordination requirements, teammate state, recent communication, and
private memory.  The collaborative system prompt describes survival, crafting,
roles, progression, and coordination.  Model outputs contain tagged fields for
an action, optional broadcast communication, and optional scratchpad state;
visible planning can also be retained in short history.

The existing reproducible baseline harness uses eight text-history turns, four
communication rounds, one previous planning turn, and one scratchpad entry.
Three model decisions are issued concurrently within a world tick.  A completed
reduced harness run used three homogeneous \texttt{gpt-5.6-luna} agents, seeds
9999--10001, all three coordination difficulties, and a 200-step cap.
\drafttodo{Confirm the public model name and immutable snapshot corresponding to
the recorded gpt-5.6-luna alias.}  This reduced run validates the baseline
pipeline and is not a market comparison or a substitute for Alem's full
20-seed, 10,000-step protocol.

\subsection{Contract Net and CU Market}

The market follows the Contract Net pattern of task announcements, bids, and
awards \citep{smith1980contractnet}.  Its implemented lifecycle is:

\begin{enumerate}
    \item publish an available mission task;
    \item collect agent bids and select an award;
    \item lease the awarded obligation for a bounded interval;
    \item verify task completion before settlement;
    \item transfer CU only after successful verification; and
    \item record failure or expiry and reauction unfinished work.
\end{enumerate}

CU is maintained in a conserved ledger backed by a 100-CU treasury.  Every
auction event is logged so balances, awards, lease transitions, settlements,
and reauctions can be replayed and audited.  The current deterministic executor
does not condition actions on CU.  Visible-CU and hidden-CU variants therefore
produce identical behavior and validate observability plumbing rather than an
incentive effect.  \drafttodo{Define the CU expansion, bid/award scoring rule,
lease duration, settlement transfer, and task-verification predicates.}

The long-term design uses auctions to distribute both tasks and role
responsibilities without a central task planner.  The current tests validate
the market infrastructure only; they do not demonstrate decentralized role
formation or role coherence.

\subsection{Completed Deterministic Market Runs}

We ran two local conditions with seed 7, eight deterministic agents, and 100
ticks.  Static direct allocation served as an engineering control, while the
market condition enabled Contract Net and CU.  Table~\ref{tab:local-market}
records the observed outcomes.

\begin{table}[t]
\centering
\small
\begin{tabular}{@{}lrrrrr@{}}
\toprule
Condition & Base & Coord. & Total & Missions & Surv. \\
\midrule
Static & 1.38\% & 0.00\% & 0.80\% & 2/6 & 8/8 \\
Market & 1.38\% & 6.29\% & 3.46\% & 2/6 & 6/8 \\
\bottomrule
\end{tabular}
\caption{Completed local engineering runs.  These conditions are not a valid
test of market efficacy because the executor and interfaces are confounded.}
\label{tab:local-market}
\end{table}

The market run produced 87 bids, 14 awards, two verified settlements, 12
reauctions, and 10 recorded failures, with exact conservation of the 100-CU
treasury.  Both conditions completed two of six mission tasks, and the market
condition lost two additional agents.  Only two of 14 awards settled.  The
increase in Total score is therefore descriptive and cannot be attributed to
the market.

\subsection{Hosted-Model Integration Check}

An eight-agent, one-tick check used \texttt{gpt-5.4-nano} through the OpenAI
Responses API.  All eight concurrent decisions returned successfully with no
provider error.  The run recorded 33,837 input tokens and 2,340 output tokens,
for 36,177 tokens total.  It establishes transport, structured generation,
action parsing, concurrency, and accounting compatibility; one tick provides
no evidence about coordination behavior.  \drafttodo{Confirm the public model
snapshot and complete request configuration for this check.}

Separately, the hosted baseline adapter records input, output, reasoning,
cached, and cache-write tokens; response status; latency; incomplete responses;
transport attempts; and normalized errors.  Stable role-specific cache keys are
sharded across concurrent traffic.  Episode-level artifacts retain resolved
configuration, source and dependency hashes, raw prompt/response traces,
actions, communication, scratchpads, trajectories, and termination reasons.

\subsection{Short-Horizon Baseline and External Context}

A pinned-upstream Alem baseline with three agents and 100 steps ended with all
three agents alive, team return 2.0, one \texttt{wake\_up} achievement, Total
score 0.266\%, and Coordination score 0\%.  This supports treating 100 ticks as
an engineering horizon rather than a meaningful Alem evaluation horizon.

For context only, the published Easy-setting results report 5.7\% Base, 6.6\%
Coordination, and 6.1\% Total for the mean of 13 LLMs; 14.1\%, 7.0\%, and
11.1\% for GPT-5.4 High; 18.1\%, 18.0\%, and 18.1\% for Gemini 3.1 Pro High;
and 21.0\%, 22.0\%, and 21.4\% for MARL trained for one billion steps
\citep{tessera2026alem}.  These values are not directly comparable to our
short runs because the published study uses three agents, language broadcasts,
scratchpad memory, much longer episodes, and multiple seeds.

\subsection{Confound Audit}

The earlier deterministic comparison cannot answer the research question for
the following reasons:

\begin{itemize}
    \item The executor navigated from full-map native state rather than the
    language agent's observation.
    \item Models received a privileged \texttt{suggested\_action} produced by
    that executor.
    \item Quote ETAs inspected hidden map, item, and coordination arrays.
    \item Market agents did not receive Alem's complete official system prompt,
    broadcast communication, or private scratchpad.
    \item History retained three compact decisions rather than eight
    observation/action turns.
    \item The market comparison used eight agents rather than Alem's canonical
    three-agent language setup.
    \item The Easy non-specialist-efficiency setting was inconsistent across
    intended and released configurations.
\end{itemize}

These mismatches affect information, action selection, memory, communication,
team size, and environment dynamics.  They preclude attributing score
differences to Contract Net or CU.

\subsection{Next Matched Pilot}

The next experiment will be a batch-size-one, paired descriptive case study.
Both treatments will use the same seed, three homogeneous LLM agents, public
model snapshot, official Alem system prompt, observation construction,
broadcast communication, scratchpad memory, history, mission, coordination
configuration, and episode horizon.  The sole treatment difference will be:

\begin{itemize}
    \item \textbf{Language only:} agents coordinate through Alem's standard
    prompt, broadcast channel, and scratchpad.
    \item \textbf{Language + market:} the same agents receive the Contract Net
    and CU interface in addition to the unchanged language condition.
\end{itemize}

The corrected market condition will remove native-state navigation,
\texttt{suggested\_action}, and hidden-array ETA access.  Market observations
and bids must be derived only from information available to the corresponding
language agent.  Both treatments will share model parameters, decoding,
timeouts, retry policy, concurrency, and accounting.

\drafttodo{Fix the paired seed, episode horizon, mission definition, CU
visibility, and exact market prompt before launching the pilot.}
With one paired seed, results will be reported as a descriptive case study with
event traces, not as a statistical efficacy claim.  Primary reporting will
include mission completion, survival, Base/Coordination/Total score, award and
settlement counts, reauctions, role/task coverage, token usage, latency, and
ledger conservation.  A later multi-seed study is required for uncertainty
estimates and any comparative claim.

\bibliography{references}

\end{document}
